\chapter{Scheduled Tasks}
\label{ch:cron}

Kiss includes a built-in task scheduler that runs Groovy scripts on a cron-like schedule.  This chapter covers the cron system, crontab configuration, and writing task scripts.

\section{The Cron System}
\label{sec:cron-system}

The \texttt{Cron}\index{Cron} class reads a schedule file and executes Groovy scripts at specified intervals.  It:

\begin{itemize}
  \item Checks for scheduled tasks every 60 seconds.
  \item Runs each task in a separate thread.
  \item Dynamically reloads the crontab file when it is modified --- no restart required.
  \item Provides a database connection to each task.
\end{itemize}

\noindent
The cron system starts automatically when the Kiss server starts, reading from \texttt{src/main/backend/CronTasks/crontab}.

\section{The Crontab File}
\label{sec:cron-crontab}

The crontab file at \texttt{src/main/backend/CronTasks/crontab} defines the schedule.  The format follows standard Unix cron conventions:

\begin{lstlisting}[style=ini]
# minute  hour  day  month  dayOfWeek  scriptName
# Run every minute
*  *  *  *  *  EveryMinute

# Run at 2 AM every day
0  2  *  *  *  NightlyCleanup

# Run at 9 AM on weekdays (Mon-Fri)
0  9  *  *  1-5  WeekdayReport

# Run at noon on the 1st and 15th
0  12  1,15  *  *  BiMonthlyExport
\end{lstlisting}

\noindent
The five time fields are:

\begin{longtable}{clP{6.5cm}}
\toprule
\textbf{Field} & \textbf{Range} & \textbf{Description} \\
\midrule
\endhead
Minute & 0--59 & Minute of the hour \\
Hour & 0--23 & Hour of the day (24-hour) \\
Day & 1--31 & Day of the month \\
Month & 1--12 & Month of the year \\
Day of week & 0--7 & 0 or 7 = Sunday, 6 = Saturday \\
\bottomrule
\end{longtable}

\noindent
Each field supports:
\begin{itemize}
  \item \texttt{*} --- matches all values.
  \item A single number --- matches that value exactly.
  \item A range --- e.g., \texttt{8-17} for 8 AM through 5 PM.
  \item A list --- e.g., \texttt{1,15} for the 1st and 15th.
\end{itemize}

\noindent
Lines starting with \texttt{\#} are comments.  The script name is the Groovy file name without the \texttt{.groovy} extension.

\section{Writing a Cron Task}
\label{sec:cron-writing}

A cron task is a Groovy file in \texttt{src/main/backend/CronTasks/}.  Every task must have a static \texttt{start(Object obj)} method:

\begin{lstlisting}[style=groovy]
import org.kissweb.database.Connection
import org.kissweb.DateTime

static void start(Object obj) {
    Connection db = (Connection) obj
    println("Task running at " +
        DateTime.currentDateTimeFormatted())

    // Perform scheduled work
    db.execute("DELETE FROM temp_data WHERE created < ?",
        DateUtils.addDays(DateUtils.today(), -30))
    db.commit()
}
\end{lstlisting}

\noindent
The \texttt{obj} parameter is a database \texttt{Connection} provided by the framework.  Key points:

\begin{itemize}
  \item The method receives its own database connection, separate from any service request.
  \item You must call \texttt{db.commit()} to persist changes.
  \item Each invocation runs in its own thread, so long-running tasks do not block other tasks or service requests.
  \item If the task throws an exception, it is caught and logged --- other tasks continue running.
\end{itemize}

\begin{warnbox}
Always call \texttt{db.commit()} in cron tasks.  Unlike service methods, the framework does \emph{not} auto-commit for cron tasks.
\end{warnbox}

\section{Example: Nightly Data Cleanup}
\label{sec:cron-cleanup}

\begin{lstlisting}[style=groovy]
// CronTasks/NightlyCleanup.groovy
import org.kissweb.database.Connection
import org.kissweb.DateUtils

static void start(Object obj) {
    Connection db = (Connection) obj

    // Delete old audit records
    int cutoff = DateUtils.addDays(DateUtils.today(), -90)
    db.execute("DELETE FROM audit_log WHERE created_date < ?",
        cutoff)

    // Deactivate expired accounts
    db.execute("""
        UPDATE users SET user_active = 'N'
        WHERE expiration_date < ? AND user_active = 'Y'
    """, DateUtils.today())

    db.commit()
    println("Nightly cleanup completed")
}
\end{lstlisting}

\noindent
Crontab entry:

\begin{lstlisting}[style=ini]
0  2  *  *  *  NightlyCleanup
\end{lstlisting}

\section{Example: Periodic Report Generation}
\label{sec:cron-report}

\begin{lstlisting}[style=groovy]
// CronTasks/WeeklyReport.groovy
import org.kissweb.database.*
import org.kissweb.*

static void start(Object obj) {
    Connection db = (Connection) obj

    Groff rpt = new Groff("WeeklyStatus-",
        "Weekly Status Report", false)
    rpt.startTable("Lw(2i) Nw(1i) Lw(1i)")
    rpt.column("Project")
    rpt.column("Hours")
    rpt.column("Status")
    rpt.endTitle()

    List<Record> projects = db.fetchAll(
        "select name, total_hours, status from project " +
        "where active = 'Y' order by name")
    for (Record r : projects) {
        rpt.column(r.getString("name"))
        rpt.column(r.getDouble("total_hours"), "C", 1)
        rpt.column(r.getString("status"))
    }

    rpt.process()
    // Report saved to temp directory
    println("Weekly report generated: " +
        rpt.getRealFileName())
}
\end{lstlisting}

\section{Dynamic Crontab Reloading}
\label{sec:cron-reload}

The cron system monitors the crontab file for changes.  If you edit the file while the server is running, the new schedule takes effect automatically on the next 60-second check cycle.  No restart is needed.

\begin{tipbox}
During development, use a short interval like \texttt{* * * * *} (every minute) to test your cron task, then change to the real schedule when it works correctly.
\end{tipbox}
